\section{Conclusion}

% In this work, we devised \textsc{TRAIL}, a new taxonomy for agentic error classification.
% We accompanied this taxonomy with an expert curated dataset of 148 agentic problem instances  consisting of 841 unique error derived from two publicly available datasets: GAIA and SWE Bench.
% We demonstrated the inability of SOTA models to act as LLM Judges, with the best performing model on our dataset (Gemini 2.5-pro) only able to achieve 18\% joint accuracy on GAIA, and 5\% joint accuracy on SWE Bench. Meanwhile, three of the eight models we test are unable to handle the complexity of TRAIL and fail to solve the task due to context-length limitations.
% The benchmark performance of the SOTA models on our dataset shows that these models are unable to systematically evaluate complex traces generated by agentic systems.
% This problem stems from the fundamental complexity of large agentic systems and finite context limitations of LLMs.
% We require a new framework in order to systematically and scalably evaluate agentic workflows.
In this work, \textsc{TRAIL}, a new taxonomy for classifying agentic errors, along with an expert-curated dataset of 148 agentic problem instances and 841 unique errors from GAIA and SWE Bench. Current SOTA models perform poorly as LLM Judges on this dataset, with \textsc{Gemini 2.5-pro} achieving only 18\% joint accuracy on GAIA and 5\% on SWE Bench; three out of eight models cannot even process the full context. These results highlight that existing models struggle to systematically evaluate complex agentic traces, due to the inherent complexity of agentic systems and LLM context limitations. A new framework is needed for scalable, systematic evaluation of agentic workflows.